% MT056 / explanation — canonical editable content.

\begin{itemize}

\item \textbf{A错误。} 深树单个模型更复杂，少量深树与大量树桩的组合能力、过拟合风险和成本不同；不能只按深度和数量断言前者一定更好。

\item \textbf{B错误。} Bagging靠多个模型之间的差异来降低波动。普通LR通常比深树稳定，换一批重采样数据后模型变化可能较小；计算便宜本身不足以证明它是很有效的bagging基学习器。

\item \textbf{C正确。} 各RF树可在自己的抽样数据上独立训练；预测时也可独立算结果，再投票或平均。因此两个阶段都能并行。

\item \textbf{D错误。} AdaBoost某轮把样本分对后，通常会降低它的相对权重，但权重仍大于0，下一轮仍参与训练。以后若它又被分错，权重还可能重新升高，所以“一次分对就永久忽略”不成立。

\item \textbf{E按原解正确。} 更多随机树使有限次平均更稳定，能降低集成的方差；但有共同错误的树无法靠数量完全相互抵消，也不保证每加一棵测试误差都单调下降。

\end{itemize}


增加树数受到的相关性限制，见\StudyRef{Answers}{mt062}{MT062}的公式。


\textbf{出处：} 2023B期中 Q5，原卷第3页（5分）


\textbf{答案依据：} 2023B期中 Q5 随卷／配套参考答案；条件与解释经复核。
